% ===== sec/05_chomsky_classification.tex =====
\section{Chomsky Classification of Agents}
\label{sec:chomsky}

We now classify trace languages by Chomsky type and exhibit, for
each type, the canonical agent shape from the empirical
\textsc{NeuroGolf ONNX ARC-AGI} system
(\texttt{experiment\_6/trace\_language.csv}).  The classification
is the literal one of Section~\ref{sec:preliminaries}; there is no
new mathematics here, only the application of an old hierarchy to a
new object.  The novelty, if any, is the structural mapping of
sub-agents in the empirical CSV to these classes, and the
architectural principle --- derived in
Section~\ref{sec:verifier-ladder} --- that the verifier's class is
the decisive design parameter.  We emphasise that the assignments
for the derived sub-agents are based on structural inspection of
their operational footprints rather than formal proofs via pumping
lemmas; the Type-0 classification of the principal generator is an
idealised upper bound (Lemma~\ref{lem:bounded-soundness}).

\subsection{Type-3: regular trace languages}

\begin{definition}[Type-3 agent]
\label{def:type3agent}
An agent $\mathcal{A}$ is \emph{Type-3} if $L(\mathcal{A}) \in
\mathcal{L}_3$, equivalently, if there exists a finite
automaton $M$ with $L(M) = L(\mathcal{A})$.
\end{definition}

In the empirical system, the Type-3 agents include the
\texttt{scanner}, \texttt{floor\_loader}, \texttt{grader}, and the
\texttt{verifier}. Their trace languages over a finite alphabet
are generated by finite-state controllers:
\begin{itemize}
\item \texttt{scanner}: the trace language is a flat sequence
of operations like \textsf{DISCOVER\_BUNDLE} read straight off
the column. No back-edge appears. No stack discipline is required.
\item \texttt{verifier}: the trace is a sequence of
\textsf{VERIFY\_TRAIN}, \textsf{VERIFY\_TEST}, etc.
The deciding finite automaton is the DFA state machine over the
pipeline stages (described formally in Section~\ref{sec:indexed}).
\end{itemize}

The Data-Driven verifier $\mathcal{V}$ of this paper is realised
as a deterministic finite automaton tracking the pipeline state.
Membership in $L(\mathcal{V})$ is decided in $O(n)$ time and $O(1)$
space per candidate trace.

\subsection{Type-2: context-free trace languages}

\begin{definition}[Type-2 agent]
\label{def:type2agent}
An agent $\mathcal{A}$ is \emph{Type-2} if $L(\mathcal{A}) \in
\mathcal{L}_2 \setminus \mathcal{L}_3$, equivalently, if there
exists a pushdown automaton $M$ with $L(M) = L(\mathcal{A})$ and
$L(\mathcal{A})$ is not regular.
\end{definition}

In the empirical system, the Type-2 sub-agents conceptually
include the \texttt{analyzer} and \texttt{miner}.  They exhibit
nested graph-building structures that naturally map to context-free
productions, such as recursively expanding nested ONNX subgraphs
(\textsf{DISCOVER\_PATTERN} followed by nested symbol resolutions).
The stack discipline is implicit in the tree-like trace of operations.

\subsection{Bounded iterative loops are strictly Type-3}

In early formulations, we tentatively linked iterative repair loops
(e.g., a \texttt{retry\_phase} loop bounded by a constant $N=3$) to
indexed grammars or cross-serial dependencies. We correct that
classification here. A loop whose iterations are bounded by a constant
$N$ does not require the unbounded stack discipline of a PDA, nor
the index stacks of an indexed grammar. It is recognizable by a
finite automaton augmented with an $\lceil \log_2(N) \rceil$-bit
counter, which is simply a larger DFA. Thus, constant-bounded
orchestration loops structurally belong to the Type-3 (regular) class.

\subsection{Type-1: context-sensitive trace languages}

\begin{definition}[Type-1 agent]
\label{def:type1agent}
An agent $\mathcal{A}$ is \emph{Type-1} if $L(\mathcal{A}) \in
\mathcal{L}_1$, equivalently, if there exists a linear-bounded
automaton $M$ with $L(M) = L(\mathcal{A})$.
\end{definition}

The Type-1 sub-agents include the \texttt{optimizer}. Its trace
language exhibits linear-bounded workspace growth: it performs
full-graph \textsf{FP16\_SURGERY}, \textsf{CAST\_COLLAPSE}, and
\textsf{DIM\_SCRUB} transformations. The length of its emitted
trace and its intermediate graph states are strictly bounded by
a constant multiple of the input graph's size, which is the LBA
hallmark.

\subsection{Type-0: recursively enumerable trace languages}

\begin{definition}[Type-0 agent]
\label{def:type0agent}
An agent $\mathcal{A}$ is \emph{Type-0} if $L(\mathcal{A}) \in
\mathcal{L}_0 \setminus \mathcal{L}_1$, equivalently, if there
exists a Turing machine $M$ with $L(M) = L(\mathcal{A})$ and
$L(\mathcal{A})$ is not context-sensitive.
\end{definition}

The Type-0 agents are the \texttt{builder} and \texttt{blender}.
The \texttt{blender} column is a prime witness: it contains a
strictly unbounded enumeration of programmatic synthesis operations
that dynamically construct and mix solvers based on arbitrary
rule criteria (\textsf{BLEND\_BUNDLE}, \textsf{Correctness Priority Rule}).
Its trace length is not bounded by a computable function of the
input length prior to encountering the verifier's bounds.

The \emph{Goal-Driven} agent of the architecture is the union of
these Type-0 sub-agents under the orchestration relation $R$ given
in Section~\ref{sec:methodology}.  Its trace language is, by
Theorem~\ref{thm:closure} run in reverse, undecidable to verify
absent the Type-3 verifier --- which is the
Robinson-MRDP~\cite{robinson1949definability,davis1961decision}
shadow of the Halting Problem in our setting.

\subsection{Summary table}

\begin{table}[ht]
\centering
\caption{Structural Chomsky classification of the empirical sub-agents in
\texttt{experiment\_6/trace\_language.csv}.}
\label{tab:chomsky-empirical}
\begin{tabular}{l|l|l}
\hline
\textbf{Sub-agent} & \textbf{Class} & \textbf{Recogniser} \\ \hline
\texttt{scanner}         & Type-3 & FA / DFA \\
\texttt{floor\_loader}   & Type-3 & FA \\
\texttt{verifier}        & Type-3 & DFA \\
\texttt{grader}          & Type-3 & FA \\
\texttt{packager}        & Type-3 & FA \\ \hline
\texttt{analyzer}        & Type-2 & PDA \\
\texttt{miner}           & Type-2 & PDA \\ \hline
\texttt{optimizer}       & Type-1 & LBA \\ \hline
\texttt{builder}         & Type-0 & TM \\
\texttt{blender}         & Type-0 & TM (emergent) \\ \hline
\end{tabular}
\end{table}

\subsection{The Goal-Driven / Data-Driven binding, formally}

The architecture of the paper is now fully formal.  Let
\[
\mathcal{G} \;=\; \mathcal{A}_{\textsf{analyzer}}
            \,\Vert\, \mathcal{A}_{\textsf{builder}}
            \,\Vert\, \mathcal{A}_{\textsf{blender}}
\]
under the orchestration relation $R_\mathcal{G}$ that allows any
of these sub-agents to consume the prefix emitted by any of the
others.  Then $L(\mathcal{G}) \in \mathcal{L}_0$.

Let
\[
\mathcal{V} \;=\; \mathcal{A}_{\textsf{verifier}}
            \,\Vert\, \mathcal{A}_{\textsf{grader}}
\]
under the orchestration relation $R_\mathcal{V}$ that requires
every \textsf{test\_*} emission to be matched against the
\textsf{tools\_index} catalogue.  Then $L(\mathcal{V}) \in
\mathcal{L}_3$.

The verified trace language is
\[
L(\mathcal{G}\,\Vert\,\mathcal{V})
\;=\; L(\mathcal{G}) \,\cap\, L(\mathcal{V})
\]
which, by Theorem~\ref{thm:closure}, lies in $\mathcal{L}_0$ and
is decidable in $O(n)$ on the verifier side.  The Type-2
sub-agents and the Type-1 sub-agents
are not explicitly hard-coded as distinct full agents inside
$\mathcal{G}$ or $\mathcal{V}$'s minimal definition; they are
\emph{projections} of the overall multi-agent system onto specific
phases of the task-solving orchestration. They emerge structurally
from the interaction between the Goal-Driven and Data-Driven bounds.

The minimality claim is therefore:
\begin{quote}
\emph{One Type-0 Goal-Driven generator and one Type-3
Data-Driven verifier suffice.  Everything else is a derivation.}
\end{quote}

This is the operational meaning of ``a single Goal-Driven agent
tamed by the Data-Driven verifier results in success.''
